\documentclass[12pt]{article}
\usepackage[a4paper,margin=2cm]{geometry}
\usepackage{amsmath,amssymb}
\usepackage{xepersian}
\settextfont[
  Path=../1401-2/HW2 - official solution/,
  Extension=.ttf,
  BoldFont=*Bd,
  ItalicFont=*It,
  BoldItalicFont=*BdIt
]{XB Niloofar}
\setdigitfont[
  Path=../1401-2/HW2 - official solution/,
  Extension=.ttf,
  BoldFont=*Bd,
  ItalicFont=*It,
  BoldItalicFont=*BdIt
]{XB Niloofar}
\setlength{\parindent}{0pt}
\setlength{\parskip}{0.55em}

\begin{document}
\begin{center}
  {\Large\textbf{اصلاحیهٔ پاسخ رسمی تمرین پنجم ـ نیم‌سال ۱۴۰۲--۲}}
\end{center}

این برگه فقط موارد نادرست پاسخ منتشرشده را اصلاح می‌کند و باید همراه فایل اصلی خوانده شود.

\section*{سؤال ۱، بخش الف}
حکم با فرض‌های نوشته‌شده در صورت سؤال درست نیست. از شرط
\(g'(x^*)=0\)
به‌تنهایی نمی‌توان نتیجه گرفت تکرار نقطهٔ ثابت دست‌کم همگرایی مرتبهٔ دو دارد. برای نمونه،
\[
  g(x)=x^*+|x-x^*|^{3/2}
\]
در \(x^*\) مشتق‌پذیر است و \(g'(x^*)=0\)، اما برای خطاها داریم
\[
  |e_{n+1}|=|e_n|^{3/2},
\]
پس مرتبهٔ همگرایی \(3/2\) است.

برای نتیجهٔ درجهٔ دو باید فرض منظم‌بودن بیشتری اضافه شود؛ برای مثال کافی است \(g'\) در همسایگی \(x^*\) لیپشیتس باشد، یا \(g''\) در آن همسایگی وجود داشته و کراندار و پیوسته باشد. تنها تحت چنین فرضی استفاده از باقی‌ماندهٔ مرتبهٔ دوم تیلور در پاسخ اصلی موجه است.

\section*{سؤال ۵}
بازهٔ انتخاب‌شده در پاسخ اصلی قابل گسترش است: مقدار مرزی
\[
  M=-\frac1b
\]
نیز معتبر است. اگر \(r\) ریشهٔ یکتای \(f\) و
\[
  x_{n+1}=x_n-\frac1b f(x_n)
\]
باشد، قضیهٔ مقدار میانگین برای عددی \(\xi_n\) میان \(x_n\) و \(r\) می‌دهد
\[
  x_{n+1}-r
  =\left(1-\frac{f'(\xi_n)}{b}\right)(x_n-r),
  \qquad
  0\leq1-\frac{f'(\xi_n)}b<1.
\]
بنابراین تکرار از ریشه عبور نمی‌کند، در همان سوی ریشه می‌ماند و به‌طور یکنوا به حدی در \([0,1]\) همگرا می‌شود. از
\[
  x_{n+1}-x_n=-\frac1b f(x_n)\longrightarrow0
\]
و پیوستگی \(f\) نتیجه می‌شود حد دنباله همان ریشه است. این استدلال حتی در حالت مجاز \(a=0\) نیز برقرار است. در پاسخ اصلی، چند نامساوی پس از ضرب در \(M<0\) با جهت نادرست نوشته شده‌اند و به حذف نابجای نقطهٔ مرزی انجامیده‌اند.
\end{document}
